\section{Manuscript saturation as a framework specialization}

The paper itself is an object produced by the research method, so the stopping rule for manuscript revision must be explicit. We define a \emph{manuscript semantic object} as a retained addition that changes at least one of the following: a scientific or method claim distinction; a citation lineage or nearest-prior-work relation; a novelty boundary; a proof or implementation obligation; an explanatory bridge needed to interpret an equation or component; a falsifier; or a reviewer-relevant evidence boundary. Sentence-level paraphrases and duplicate citations do not count as semantic growth.

For manuscript pass \(r\), let
\[
\mathcal M_r=\operatorname{Canon}
(\text{retained manuscript semantic objects through pass }r),
\qquad
\Delta_r^{\mathrm{paper}}=\mathcal M_r\setminus\mathcal M_{r-1}.
\]
A pass is semantically flat when \(\Delta_r^{\mathrm{paper}}=\varnothing\) after deduplication and it creates no new unresolved proof obligation, novelty correction or reviewer-blocking contradiction. Flatness is recorded separately for the registered lenses: epistemology/formal methods, mathematical geometry, scientific-agent literature, and methods/reproducibility.

A \emph{locally saturated publication projection} additionally requires: all mandatory search/review routes are covered or explicitly blocked; the freshness scan reaches its registered cutoff; nearest-work and citation-equivalence audits are complete for the main novelty claims; structural proof obligations are resolved or explicitly open; and every remaining manuscript issue is classified as empirical/deferred, out of scope, or blocked by unavailable evidence. A material open issue prevents the certificate. Page count never supplies saturation credit.

This protocol is intentionally a specialization of the existing saturation/stopping and review surfaces, not a new 25th high-impact research-method surface. The underlying mechanism is the same semantic-growth accounting already used for research fibers. What changes is the object being measured and the terminal classification of open items. Same-context manuscript review can support only \emph{local manuscript flatness}. It cannot be relabelled as independent peer review or independent scientific saturation.

The reopen rule is strong. Any later exogenous concept, citation, counterexample, implementation result or reviewer comment that materially changes the novelty boundary, formal architecture or evidence interpretation adds a new manuscript semantic object and reopens the projection. This makes the paper a versioned interface to a living atlas rather than a claim that the literature has been permanently exhausted.

The protocol also disciplines length. In response to a request for ``more depth,'' the correct action is not to multiply prose uniformly. A section grows when the reader currently lacks an intellectual ancestor, formal distinction, proof obligation, empirical boundary or explanatory bridge needed to evaluate the claim. A paragraph that restates an equation without adding such information is removed. Saturation therefore opposes both premature compression and decorative expansion.

\subsection{Section purpose is itself a proof obligation}

A saturation protocol needs a deletion rule as well as an addition rule. For every top-level section, the manuscript records a \emph{reader obligation}: the question that section must answer which is not already answered elsewhere. If no such obligation can be named, the section should be merged or removed even when its prose is correct. This turns the editorial question ``why is this page necessary?'' into a structural audit rather than a stylistic preference.

The present paper uses the following section-purpose map. The Introduction establishes the research problem and the claim boundary. Intellectual lineage establishes what is inherited and therefore what novelty remains. Formal method defines the scientific state and transition semantics required to evaluate the claims. Scientific memory explains how raw evidence becomes that state rather than an ordinary RAG store. The atomic lifecycle translates the formalism into an auditable execution contract. The Obsidian--J-space--RAKL comparison prevents category mistakes among navigation, representation and epistemic geometries. Workspace and OWMD explain bounded cognition and ontology escape. The known-answer trace tests deterministic mechanics. Bounded scientific cognition states the resource boundary. Method learning and Self-RAKL state how the method may change without self-authorization. Reference-implementation correspondence binds formal claims to code. The preregistered evaluation programme defines what empirical evidence would justify stronger claims. Manuscript saturation defines the stopping rule for this publication projection. Discussion exposes rival interpretations and failure modes. Reproducibility identifies the artifact and what can be replayed. Structural proof obligations make the mathematical words falsifiable.

This map supplies a concrete redundancy test. If a paragraph in Discussion only restates a formal definition, it should be removed unless it interprets a consequence or boundary. If Related Work merely names systems already used to narrow a specific contribution, it should be reorganized around the novelty claim rather than expanded as a catalogue. If the known-answer demo repeats the lifecycle description without showing a selective failure or invariant, it adds no evidential role. Conversely, a short section may need expansion when its reader obligation cannot yet be discharged. The three-geometry section grew for this reason: the earlier analogy to Obsidian did not establish how a note graph, a cone-like representation space and a typed epistemic state differ mathematically.

The same rule applies below the section level. A paragraph should normally contribute one of a small number of functions: define an object, locate intellectual ancestry, state evidence, derive a consequence, distinguish nearby concepts, expose a limitation, specify a falsifier, or connect a formal object to implementation/evaluation. A paragraph whose only function is rhetorical emphasis is a candidate for deletion. Equations and tables receive the same treatment. An equation is justified when its symbols or invariants matter to later reasoning; a table is justified when simultaneous comparison prevents ambiguity that prose would hide. Mathematical decoration and citation decoration are both failures of purpose.

This purpose audit is deliberately coupled to depth. ``Concise'' does not mean omitting the argument required to assess a claim, and ``comprehensive'' does not mean saying the same thing three ways. The target is \emph{epistemic density}: retained words should increase the reader's ability to locate the claim, understand its ancestry and assumptions, reproduce its derivation, test its boundary, or identify the evidence needed to change it. Length may therefore increase sharply during an early saturation pass and then decrease during canonicalization as duplicate explanations are merged.

The page budget creates a projection problem rather than a truth problem. If a journal format cannot contain all canonical detail, the solution is to move reproducible derivations, extended lineage matrices or machine-readable receipts to appendices and repository artifacts while preserving enough in the main paper to make every central transition intelligible. The omitted material remains addressable through stable pointers. Silent deletion of a blocking assumption or counterexample to meet a page limit is not acceptable compression; it changes the epistemic content of the publication projection.

A final editorial pass therefore asks two questions in both directions. For every retained page: what reader obligation would become materially harder to discharge if this content disappeared? For every central claim: is there a page that actually discharges the corresponding ancestry, formal, evidence and boundary obligations? A ``no'' to the first question triggers deletion or merging. A ``no'' to the second triggers a new manuscript fiber. The paper stops changing only when these operations, together with the external literature and formal audits, become semantically flat. The criterion is intentionally symmetric: it prevents ornamental pages while reopening underdeveloped claims whose evidence, ancestry or proof obligations remain hidden.

\subsection{Operational saturation ledger}

A manuscript-saturation run maintains a ledger rather than a subjective sense of completion. Each review pass records its lens, route, context identifier, sources inspected, semantic objects proposed, semantic objects retained after canonicalization, and open items created or closed. A retained object receives a stable identifier so later passes can distinguish new content from repeated rediscovery.

The ledger separates seven semantic kinds: claim distinction, citation cluster, novelty-boundary change, proof obligation, explanatory bridge, falsifier and review repair. This taxonomy is intentionally coarse. Its purpose is to answer whether the paper's epistemic content grew, not to grade prose. A new citation that supports an already well-established sentence without changing lineage may be useful copyediting but adds no semantic object. A single prior paper that shows a claimed contribution is established can create a major novelty-boundary object even though it adds only one reference.

Flatness is computed after canonicalization. If three review lenses independently rediscover the same missing distinction, the first retained instance is growth and the others are corroboration of its importance. Conversely, splitting one idea into many headings does not manufacture novelty. This makes semantic saturation resistant to cosmetic expansion.

Open items have terminal classifications. `MATERIAL\_OPEN' blocks local saturation because it could change the paper's central claims or architecture. `EMPIRICAL\_DEFERRED' identifies a claim that can only be resolved by the preregistered experiment or other future evidence. `OUT\_OF\_SCOPE' records a neighboring question whose exclusion is deliberate and justified. `BLOCKED\_MISSING\_EVIDENCE' records a relevant question that cannot currently be checked. These states are visible in the saturation receipt so the absence of prose is not misread as an implicit solution.

A route or lens can be flat while the whole manuscript remains active. For example, a mathematics pass may add no new object while a literature freshness pass discovers a new scientific-agent benchmark that changes the evaluation section. That new object resets the global flat tail and can reopen earlier sections because its implications propagate. The stopping rule therefore operates on the manuscript dependency graph, not on a fixed number of passes per section.

Saturation also has a cost boundary. The process is not required to read every paper ever written. Instead it must show that the registered routes reached local semantic flatness, that nearest competing work was explicitly sought, and that unresolved omissions are classified. A future external reviewer can then challenge the certificate constructively by naming the missing semantic object. If the challenge is material, the proper response is to reopen the paper, not to defend the word ``saturated.''

In the present release, the four review lenses were chosen because they correspond to the most consequential failure families: epistemology/formal methods can expose authority or identification errors; mathematical geometry can expose category mistakes in lattice, graph, cone and sheaf language; scientific-agent literature can expose novelty and evaluation gaps; and methods/reproducibility can expose untraceable claims or artifact drift. Other journals or domains can register additional lenses without changing the underlying saturation mechanism.

\subsection{What a manuscript-saturation certificate actually certifies}

A saturation certificate must identify its universe. For a manuscript projection \(P\), define a registered audit scope
\[
\Sigma_P=(Q_P,L_P,R_P,D_P,t^\star,B_P),
\]
where \(Q_P\) is the set of central paper claims and proof obligations, \(L_P\) the review lenses, \(R_P\) the search-route families, \(D_P\) the declared source universe or source classes, \(t^\star\) the freshness cutoff, and \(B_P\) the resource and page/context budget. The certificate says that semantic growth became flat under \(\Sigma_P\) after the registered audits. It says nothing about sources, disciplines or future evidence outside that tuple.

The route set matters because repeated searching with one vocabulary can create false flatness. A claim may appear saturated under exact-term search while remaining incomplete under historical terminology, mathematical equivalence, adjacent-domain function search or citation-neighborhood expansion. The paper therefore inherits the OWMD lesson: at least one route for each novelty-critical mechanism should be capable of succeeding without repeating the framework's own labels. The Obsidian miss is the motivating counterexample. The concept was easy to recognize after it was supplied, yet the earlier route family had not surfaced it.

The claim set matters for a different reason. A manuscript can be globally long and still be locally thin at its most important inferential bridge. Saturation is not uniform word density. The audit asks whether every central claim has enough ancestry, formal semantics, evidence boundary and falsifier for a reviewer to determine what would make it wrong. A background paragraph can remain short because its only job is orientation. A geometry section may need substantial depth because a category mistake there would invalidate several downstream claims. The page budget therefore constrains presentation, not the amount of canonical research state that must exist behind the paper.

Proof obligations are treated as semantic objects because a newly discovered obligation changes what the paper must establish even before the proof is supplied. For example, using the word ``lattice'' creates meet/join or closure-system obligations; using ``sheaf'' creates locality, restriction and gluing obligations; claiming independent corroboration creates evidence-lineage obligations; claiming a matched experiment creates treatment-definition and resource-comparability obligations. A pass that discovers such an obligation is non-flat even if it adds no new empirical result. The manuscript remains active until the obligation is discharged, explicitly weakened, or classified open.

Citation saturation is similarly semantic rather than numeric. One hundred references that all repeat the same conceptual neighborhood may add less coverage than one remote paper that changes the novelty boundary. Each source is therefore mapped to a role: foundational ancestor, nearest competing method, mathematical precedent, measurement/evaluation precedent, counterexample, contemporary benchmark, implementation analogue or scope boundary. The useful question is not ``does every paragraph contain citations?'' but ``could a knowledgeable reviewer name a materially closer ancestor or counterexample that the current claim map does not represent?'' When such a source appears, the novelty coordinate reopens.

The certificate also separates \emph{flatness} from \emph{closure}. A pass can add no new semantic object because an important database is unavailable. That is flat observation, not closure. Likewise, a route can be exhausted because its token budget was reached while new concepts were still arriving. The correct states are blocked or active-non-flat. Closure requires that material open items are absent and that deferred empirical questions are explicitly labeled. This prevents resource exhaustion from masquerading as knowledge saturation.

Same-context recursion introduces a final dependence problem. The same researcher or model can run multiple lenses and still share hidden assumptions, search habits and evidence ancestry. Repeated flat passes are useful engineering evidence that the current context has stopped generating new repairs, but they cannot establish independent literature completeness or peer-review consensus. The receipt therefore has two separate fields: same-context local manuscript saturation and independent review status. In this release the former can become true while the latter remains false.

A practical certificate should be falsifiable by a small object. An external reviewer does not need to refute the entire paper to reopen it; supplying one materially new prior method, one counterexample to a formal claim, one missing proof obligation, one unrepresented failure mode or one artifact mismatch is enough. The challenge is canonicalized into the ledger, the affected dependency neighborhood reopens, and later flatness must be re-earned. In this sense saturation is a versioned state, not a medal attached permanently to the manuscript.

This interpretation makes the paper itself an executable example of epistemic mechanics. The canonical project atlas can be richer than the publication, while the publication declares which slice it projects and which open coordinates it intentionally leaves out. A reader can therefore distinguish three statements that conventional writing often conflates: ``the authors stopped editing'', ``the registered review process stopped producing new material objects'', and ``the scientific field contains no relevant unknowns.'' Only the second is the manuscript-saturation claim made here.
